\chapter{Conception}
\label{chap:conception}

\section{Introduction}
% Annoncer le contenu : architecture generale, conception detaillee
% (diagrammes de sequence), modele de donnees, conception de l'IHM.

[\`A compl\'eter]

\section{Architecture g\'enérale}
% Decrire l'architecture globale de la solution :
%   - Architecture 3-tiers : Frontend (Angular) <-> Backend (API REST,
%     port 8081) <-> Bases de donnees / moteur Drools
%   - Keycloak comme serveur d'authentification (port 8080)
%   - Schema d'architecture a inserer (cf. DOCUMENTATION_IA_NLP.md
%     section 12 pour un exemple de schema en couches pour le module IA)

La plateforme E-FORCE repose sur une architecture en trois couches~:

\begin{itemize}
    \item \textbf{Couche pr\'esentation} : application Angular 21
    (standalone components), responsable de l'interface utilisateur et
    de la navigation ;
    \item \textbf{Couche m\'etier / services} : API REST Java (Spring
    Boot, port 8081) expose les ressources (utilisateurs, processus,
    t\^aches, r\`egles, dossiers import/export, NLP) et int\`egre le moteur
    de r\`egles \textbf{Drools} ;
    \item \textbf{Couche s\'ecurit\'e} : serveur \textbf{Keycloak} (port
    8080) g\'erant l'authentification SSO et les r\^oles des utilisateurs.
\end{itemize}

[\`A compl\'eter : ins\'erer le sch\'ema d'architecture globale.]

% \begin{figure}[H]
%     \centering
%     \includegraphics[width=\textwidth]{images/architecture_globale.png}
%     \caption{Architecture g\'en\'erale de la plateforme E-FORCE}
%     \label{fig:architecture-globale}
% \end{figure}

\section{Architecture du frontend Angular}
% Decrire l'organisation des dossiers du projet Angular (cf.
% DOCUMENTATION.md section 3) : core (services, guards, interceptors),
% components (regle, bpmn-viewer, nlp-input, drl-preview), features
% (dashboard, processus, users, taches), pages (import/export).

[\`A compl\'eter : pr\'esenter l'organisation des dossiers \texttt{core/},
\texttt{components/}, \texttt{features/} et \texttt{pages/} du projet
Angular, \'eventuellement sous forme d'arborescence ou de diagramme de
paquetages.]

\section{Conception d\'etaill\'ee}
% Pour chaque cas d'utilisation cle, fournir un diagramme de sequence
% illustrant les echanges entre l'utilisateur, le frontend Angular,
% les services et le backend. Exemples suggeres :
%   - Authentification (login via Keycloak)
%   - Creation d'une regle metier via l'IA (NLP)
%   - Generation et visualisation d'un diagramme BPMN

\subsection{Diagramme de s\'equence : authentification}
[\`A compl\'eter : diagramme de s\'equence du processus de connexion via
Keycloak.]

\subsection{Diagramme de s\'equence : cr\'eation d'une r\`egle via l'IA (NLP)}
% Le pipeline est deja documente dans DOCUMENTATION_IA_NLP.md section 9
% (flux utilisateur complet) : Saisie -> POST /api/nlp/convertir ->
% Apercu (confiance, conditions, DRL) -> Edition eventuelle -> POST
% /api/regles

[\`A compl\'eter : diagramme de s\'equence illustrant le flux Saisie ->
Conversion IA -> Aper\c{c}u -> (\'Edition) -> Cr\'eation de la r\`egle.]

\subsection{Diagramme de s\'equence : g\'en\'eration du diagramme BPMN}
[\`A compl\'eter : diagramme de s\'equence illustrant la g\'en\'eration du
XML BPMN \`a partir des t\^aches d'un processus et son affichage dans le
visualiseur (\texttt{bpmn-js}).]

\section{Mod\`ele de donn\'ees}
% Decrire les principales entites manipulees par l'application, sous
% forme de diagramme de classes UML. Les modeles TypeScript existants
% (src/app/models/) donnent une bonne base :
%   - Utilisateur (Keycloak)
%   - Processus / Tache
%   - RegleMetier / Condition / Categorie / Version
%   - DossierImport / DossierExport
%   - NlpConversionRequest / NlpConversionResult / NlpRegleDto

[\`A compl\'eter : diagramme de classes (ou mod\`ele relationnel) reprenant
les entit\'es \texttt{Processus}, \texttt{Tache}, \texttt{RegleMetier},
\texttt{Condition}, \texttt{Categorie}, \texttt{DossierImport},
\texttt{DossierExport}, etc., et leurs relations.]

% \begin{figure}[H]
%     \centering
%     \includegraphics[width=\textwidth]{images/diagramme_classes.png}
%     \caption{Diagramme de classes (extrait)}
%     \label{fig:diagramme-classes}
% \end{figure}

\section{Conception de l'interface utilisateur}
% Presenter les maquettes (wireframes) des principaux ecrans :
% dashboard, liste des regles, formulaire de creation via IA, liste des
% processus, visualiseur BPMN, import/export.

[\`A compl\'eter : maquettes ou \'ecrans repr\'esentatifs de l'application.]

\section{Conclusion}
% Resumer le chapitre et faire la transition vers la realisation.

[\`A compl\'eter]

\clearpage
